% SPDX-License-Identifier: Apache-2.0
\section{A Flash Read As Memory}
\label{sec:x-flashwin}

A program lives in the instruction memory's initial contents, so
changing it means building a bitstream. The board's configuration
flash has room beside the bitstream for programs, and the master in
Section~\ref{sec:x-spi} can read it, a byte at a time, under software
that knows the command. That is the wrong shape for a bootloader and
the wrong shape entirely for a core that fetches from the bus, because
neither can run software to read the software it is about to run.

\code{FlashWin} is the other way round: an address range in which an
ordinary read is answered with what the flash holds there. Nothing is
sequenced by software, and nothing has to be. A read at an address
sends the fast read command \code{0x0b}, the three bytes of that
address with its low two bits cleared, one byte the chip throws away,
and then takes four, which are the word, lowest address in the lowest
byte.

Nine bytes on the wires for four on the bus is the trade. At the
divider the run below uses, a word takes about three hundred cycles,
which is slow beside a memory and fast beside a resynthesis. A cache
in front of the window is the answer if it matters, and there is none
here.

The wires are mode 0, and only mode 0: the clock idles low, the chip
takes a bit as the clock rises, and the master moves one as it falls.
That is the mode every flash answers a fast read in, so the window
needs neither the two mode bits the master has nor the register to
hold them. The select is held low for the whole of a read and released
after it, so each read is a command the chip sees whole, which is what
a real chip requires and what the model checks.

A write is answered \code{SLVERR} rather than ignored. The window is
read only, and saying so is what makes it safe to fetch from: nothing
on the bus can change what the instruction stream is reading, so the
question of what a fetch sees when a write lands in the same cycle
does not arise.

The run reads two words that the chip model holds, checks each against
what was put there, and then writes, which comes back refused. The
chip is \code{FlashDevice}, the same model that checks the master, so
the two ways of reaching a flash are checked against one description
of one chip.

\begin{figure}[htbp]
\centering
\input{flashwin_timing.tex}
\caption{The start of a read: the select goes low, and the command
goes out on \code{mosi} a bit at a time as \code{sclk} turns, with
\code{phase} counting the nine bytes.}
\label{fig:x-flashwin}
\end{figure}

\lstinputlisting[caption={\code{ex\_flashwin.rs}. Run with \code{bazel run //lib/examples:ex\_flashwin}.},label={lst:x-flashwin}]{lib/examples/ex_flashwin.rs}

\lstinputlisting[style=ascii,caption={What \code{ex\_flashwin} printed when the build ran it: the two words read out of the chip, and the write it refused.},label={lst:x-flashwin-out}]{out_flashwin.txt}
